% Paper 3A Supplementary Information. Built and checked by paper/tex/build_ms.py with main.tex.
\documentclass[11pt]{article}
\usepackage[a4paper,margin=25mm]{geometry}
\usepackage[T1]{fontenc}
\usepackage{lmodern}
\usepackage{amsmath,amssymb}
\usepackage{graphicx}
\usepackage{framed}
\usepackage{booktabs}
\usepackage{textcomp}
\usepackage[hidelinks]{hyperref}
% Macros shared by main.tex and si.tex. The build gates (paper/tex/build_ms.py) read these names.
% \V{key}{value}  a number (math mode); key -> paper/tex/bindings.json; G1 compares value with the source
% \W{key}{word}   a count written as a word (text); same binding rules (fmt word / Word)
% \C{...}         a mathematical constant inside a formula (2 in 2 pi, 0 in p -> 0); no data
% \ID{...}        an identifier that contains digits (hash, date, version, DOI, ORCID, model-cell names); no data
% \snote{n}       a Supplementary Note number; G2 checks that si.tex defines it
% \pending{name}{text}  a result not yet available; G7 lists it and a --final build refuses it
\newcommand{\V}[2]{\ensuremath{#2}}
\newcommand{\W}[2]{#2}
\newcommand{\C}[1]{#1}
\newcommand{\ID}[1]{#1}
\newcommand{\snote}[1]{#1}
\newcommand{\pending}[2]{\par\noindent\fbox{\parbox{0.95\linewidth}{\textit{Pending result (#1).} #2}}\par}
\newcommand{\pss}{p/p_{\mathrm{sat}}}
\newcommand{\qref}{q_{\mathrm{ref}}}
\newcommand{\Qst}{Q_{\mathrm{st}}}


\title{Supplementary Information for: Hard thermodynamic constraints and the temperature extrapolation of learned
water isotherms in metal-organic and covalent organic frameworks: a pre-registered test}
\author{Abdulhamid Batayhi}
\date{}

\begin{document}
\maketitle

\section*{Supplementary Note~\snote{1}: Can each verdict fail? Synthetic worlds and power}

Before the freeze, the whole pipeline (the tuned configurations, fitting from \V{n_seeds}{5} seeds, the violation
counts and the analysis code) was run in synthetic worlds that share no uptake data with the confirmatory set. Each
world has \V{n_conf}{13} materials; material $i$ uses the temperatures of confirmatory material $i$. An expected
outcome was written for each world before it ran. In KS0 the truth is a two-site Sips isotherm with van't Hoff
affinities, inside the family of P; in DA it is a Dubinin-Astakhov curve, inside the family of the hard
characteristic curve (A-hard); in FW it is a Kirchhoff-Sips isotherm with a large heat capacity, so that van't
Hoff is wrong; FLAT has one weakly curved Langmuir site; SUB has a sub-linear Sips isotherm outside the family of P
and inside that of the penalty model V-pen.

The comparison of P with A-hard discriminated in both directions: P extrapolated better in KS0
($R = \V{ks0_h1c}{0.52}$, exact \V{ci_main}{95}~\% interval \V{ks0_h1c_lo}{0.35} to \V{ks0_h1c_hi}{0.70}) and
A-hard better in DA ($R = \V{da_h1c}{1.26}$, \V{da_h1c_lo}{1.07} to \V{da_h1c_hi}{1.49}). \W{n_missed_word}{Three}
expectations were missed. In FW the comparison of P with the free temperature network F-hard still read ``P
better'' ($R = \V{fw_h1a}{0.33}$): with two training temperatures the free network cannot extrapolate, so this
comparison cannot fail and is reported descriptively. In FLAT the unconstrained network F-none violated
physics ``often''. In SUB, built for V-pen to win, P was slightly better ($R = \V{sub_h1b}{0.96}$,
\V{sub_h1b_lo}{0.94} to \V{sub_h1b_hi}{0.99}); the mixture of Langmuir terms in P approximates a sub-linear isotherm
over the measured range. The accuracy comparison of hard constraints with penalties therefore could not be made to
fail synthetically, and it is tested as non-inferiority.

The power of each verdict was simulated with the frozen analysis code at $n = \V{n_conf}{13}$
(\V{mde_nsim}{2000} simulations per setting). At the spread expected for the comparison with V-pen
(\V{h1b_sigma}{0.158} decades), the non-inferiority criterion is met with probability \V{ni_size_sim}{0.0225} when
the true ratio is at the margin, $R = \V{sesoi_hi}{1.25}$, and \V{ni_power_n}{0.51} (normal) or
\V{ni_power_t}{0.74} (Student $t$, \V{t_df}{3} degrees of freedom) when the two models are equally accurate. A
``not met'' outcome is therefore likely even under equal accuracy and is never read as evidence that the
constraints cost accuracy.

% written by build_ms.py from refs_master.json (order of first citation); do not edit
\renewcommand\refname{References}
\begin{thebibliography}{99}
\bibitem{oppenheim2025} Oppenheim, J.; Yang, Z.; Dinakar, B.; Dinc\u{a}, M. High-capacity water sorbent cycles without hysteresis under dry conditions. \textit{Nat. Commun.} \textbf{2025}, \textit{16} (1), 4297. DOI: 10.1038/s41467-025-59551-2.
\bibitem{kang2024} Kang, M. S.; Heo, I.; Park, S. H.; Bae, J.; Kim, S.; Kim, G.; Kim, B.-H.; Jeong, N. C.; Yoo, W. C. Time-efficient atmospheric water harvesting using fluorophenyl oligomer incorporated MOFs. \textit{Nat. Commun.} \textbf{2024}, \textit{15} (1), 9793. DOI: 10.1038/s41467-024-53853-7.
\bibitem{liu2022emm8} Liu, Z.; Xu, J.; Xu, M.; Huang, C.; Wang, R.; Li, T.; Huai, X. Ultralow-temperature-driven water-based sorption refrigeration enabled by low-cost zeolite-like porous aluminophosphate. \textit{Nat. Commun.} \textbf{2022}, \textit{13} (1), 193. DOI: 10.1038/s41467-021-27883-4.
\bibitem{gao2026} Gao, M.-Y.; Bezrukov, A. A.; Eaby, A. C.; Deng, C.; Liu, L.; Song, B.-Q.; Nikkhah, S. J.; Vandichel, M.; Zaworotko, M. J. Molecular-level insight into water adsorption and projected atmospheric water harvesting performance in a hydrolytically stable MOF. \textit{Nat. Commun.} \textbf{2026}, \textit{17} (1), 9494. DOI: 10.1038/s41467-026-76301-0.
\bibitem{wagner1993} Wagner, W.; Pruss, A. International equations for the saturation properties of ordinary water substance. Revised according to the International Temperature Scale of 1990. \textit{J. Phys. Chem. Ref. Data} \textbf{1993}, \textit{22} (3), 783--787. DOI: 10.1063/1.555926.
\end{thebibliography}


\end{document}
